\documentclass[11pt,a4paper]{article}
\usepackage{labreport}

\LabCourse{@@COURSE@@}
\LabCourseShort{@@COURSESHORT@@}
\LabReportType{@@REPORTTYPE@@}
\LabNumber{@@LABNUMBER@@}
\LabTitle{@@TITLE@@}
\LabStudent{@@STUDENT@@}
\LabStudentID{@@STUDENTID@@}
\LabMajor{@@MAJOR@@}
\LabCollege{@@COLLEGE@@}
\LabDate{@@DATE@@}
\LabWordmarkFile{assets/university_wordmark.png}
\LabSealFile{assets/university_logo.jpeg}
\LabHeaderLeft{@@HEADERLEFT@@}
\LabHeaderRight{@@HEADERRIGHT@@}

\begin{document}

\LabMakeCover
\LabBeginMain

\begin{abstract}
\noindent
\LabPending{Write the abstract after the measured results and exclusions are known.}
\end{abstract}

\section{Introduction and Scope}

State the experimental question, the system under test, and the claims the report will support.

\section{Experimental Environment and Method}

Describe the equipment, software, topology or apparatus, variables, controls, and measurement procedure. Include the addresses, ports, parameters, and configuration values needed to reproduce the run.

\section{Results}

Present the measured observations before interpretation. Use tables for multi-condition comparisons and figures for apparatus, topology, or captured evidence.

\section{Discussion}

Connect each result to the mechanism under test. Separate strong causal evidence, correlational observations, and mechanistic hypotheses.

\section{Conclusion}

\LabPending{State the accepted result, the measured boundary conditions, and the remaining unresolved question.}

\appendix
\section{Evidence Checklist}

List the raw observations, configuration exports, screenshots, and generated data needed to reproduce each claim.

\begin{thebibliography}{9}
\footnotesize
\bibitem{course}
Course laboratory sheet or apparatus documentation.
\end{thebibliography}

\end{document}

